% !TeX program = pdflatex
\documentclass[10pt,letterpaper]{article}
% Compiler Engineering Suite - shared standalone pdfLaTeX preamble
\documentclass[10pt,letterpaper]{article}
\usepackage[T1]{fontenc}
\usepackage[utf8]{inputenc}
\usepackage{lmodern}
\usepackage[margin=.78in,headheight=18pt]{geometry}
\usepackage[table]{xcolor}
\usepackage{microtype}
\usepackage{amsmath,amssymb}
\usepackage{booktabs,longtable,array}
\usepackage{enumitem}
\usepackage{fancyhdr}
\usepackage{hyperref}
\usepackage{needspace}
\usepackage{fancyvrb}
\usepackage{tikz}
\usetikzlibrary{arrows.meta,positioning,shapes.geometric}
\definecolor{CompilerBlue}{HTML}{1F4E79}
\definecolor{CompilerInk}{HTML}{263746}
\definecolor{CompilerGray}{HTML}{52616B}
\definecolor{CompilerLight}{HTML}{EFF4F8}
\hypersetup{colorlinks=true,linkcolor=CompilerBlue,urlcolor=CompilerBlue,citecolor=CompilerBlue}
\setlength{\parindent}{0pt}
\setlength{\parskip}{6pt plus 2pt}
\setlist[itemize]{leftmargin=*,itemsep=3pt,topsep=2pt}
\setlist[enumerate]{leftmargin=*,itemsep=3pt,topsep=2pt}
\pagestyle{fancy}
\fancyhf{}
\fancyhead[L]{\small\textcolor{CompilerGray}{Compiler Engineering Documentation Suite}}
\fancyhead[R]{\small\thepage}
\renewcommand{\headrulewidth}{.4pt}
\newcommand{\evidence}[1]{\par\smallskip\noindent\textbf{Evidence required:} #1\par}
\newcommand{\sourceanchor}[1]{\textbf{Book anchor:} #1\par}
\newcommand{\modern}{\textbf{Modern engineering extension (not attributed to the source book):} }
\newcommand{\writinglines}[1]{\par\noindent\rule{\linewidth}{.3pt}\par\vspace{#1\baselineskip}\noindent\rule{\linewidth}{.3pt}\par}
\newcommand{\checkboxline}[1]{\noindent $\square$\quad #1\par}
\newcommand{\formfield}[2]{\noindent\textbf{#1}\par\vspace{#2\baselineskip}\hrule\medskip}

\hypersetup{pdftitle={Formal Methods and Analysis Reference},pdfauthor={Compiler Engineering Documentation Suite}}
\begin{document}
\begin{titlepage}
\centering\vspace*{1.0in}
{\Huge\bfseries\textcolor{CompilerBlue}{Formal Methods and Analysis Reference}\par}
\vspace{.22in}{\Large Finite automata, CFG fixed points and legality proofs\par}
\vspace{.4in}{\large Compiler Engineering Documentation Suite\par}
\vfill
\begin{minipage}{.86\linewidth}\small
Source basis: Keith D. Cooper and Linda Torczon, \emph{Engineering a Compiler}, Third Edition (Morgan Kaufmann/Elsevier, 2023). All explanatory text, examples, questions, and implementation designs in this document are independently authored. The book is not included.\par\medskip
Document class: Technical reference.
\end{minipage}
\vfill\textcolor{CompilerGray}{Version 1.0 --- 2026-10-30}\par
\end{titlepage}
\tableofcontents\clearpage
\section{Formal model checklist}
\begin{itemize}
\item Specify alphabet, grammar, lexical priorities, tree operator semantics, and source-language types.
\item Define abstract state domain and order, top/bottom elements, meet/join, transfer and boundary values.
\item For each rewrite, state preconditions, invariant preservation and semantic equivalence.
\item Treat compiler crash freedom and language acceptance as different correctness properties.
\end{itemize}

\section{Selected formal equations}
For backward liveness on a CFG with successors $succ(b)$:
\[
 OUT[b]=\bigcup_{s\in succ(b)}IN[s],\qquad
 IN[b]=USE[b]\cup(OUT[b]\setminus DEF[b]).
\]
For dominance, entry $d(entry)=\{entry\}$; other blocks initialize to all blocks and iterate
\[
 Dom(b)=\{b\}\cup\bigcap_{p\in pred(b)}Dom(p).
\]
For abstract interpretation, a monotone function $F:L\to L$ is iterated to a fixed point over an appropriate finite-height domain. Ensure termination conditions and widening assumptions are documented for infinite domains.

\section{Original proof obligation register}
\begin{longtable}{|p{0.20\linewidth}|p{0.38\linewidth}|p{0.35\linewidth}|}
\hline
Claim & Assumption & Discharge method\\ \hline\endfirsthead
\hline
Claim & Assumption & Discharge method\\ \hline\endhead
Constant fold & Operator has defined pure semantics & Compare evaluator results over boundary-domain samples\\ \hline
DCE & Removed instruction has no observable effects & Live-out and effect-trace proof\\ \hline
Code motion & Dominance, alias and exception safety hold & Path-sensitive equivalence tests\\ \hline
Scheduling & Dependences and resources respected & Trace equivalence + cycle simulation\\ \hline
Register allocation & Interference graph conservatively models overlap & No shared color on adjacent nodes\\ \hline
\end{longtable}

\section{Worked example: liveness}
Suppose $b_1$ defines $x$ then branches to $b_2$ and $b_3$. Both successors use $x$ before redefining it. Consequently $x\in OUT[b_1]$, but whether it belongs to $IN[b_1]$ depends on whether $USE[b_1]$ contains $x$ before its definition. This distinction prevents the common error of equating ``live-out'' with ``live-in.

\end{document}
